\section{Discussion}
To further highlight the advantages of our framework over existing methods~\cite{peng2024d,huang2025real}, we discuss it from three perspectives: deployment efficiency, scalability, and training efficiency.

\noindent \textbf{Deployment Efficiency.}
Our method introduces zero modification to detector architectures and does not alter the inference pipeline, ensuring that no additional computational cost or latency is incurred. This deployment-friendly property is crucial for industrial applications, where real-time detectors are tightly integrated into existing systems and hardware-constrained environments.

\noindent \textbf{Scalability.}
The framework is highly general and can be seamlessly applied to detectors with diverse architectures, including CNN-based and transformer-based detectors. It enables all types of real-time detectors to quickly benefit from the rapid progress of Vision Foundation Models (VFMs). 
Moreover, the framework is not restricted to any specific type or scale of VFM.
It can flexibly incorporate different foundation models, such as DINOv3~\cite{simeoni2025dinov3}, MAE~\cite{he2022masked}, or CLIP~\cite{radford2021learning}, and even benefit from arbitrarily large models for distilling semantics into real-time detectors.
This flexibility also opens up promising directions for multi-VFM semantic integration, further demonstrating the framework's generality and scalability.

\noindent \textbf{Training Efficiency.}
Our approach is lightweight and easy to implement. Since neither the detector structure nor the optimization pipeline is modified by the incorporation of VFMs, the additional training cost remains minimal. This efficiency highlights the practicality of our method for large-scale applications and real-time industrial deployment.